\chapter{Stellar Astrophysics}

\section*{Theory problems}

\probhead{7.1}{}{2007}{Chiang Mai, Thailand}{TH 1.9}{2}
% source id: 2007_TH_1.9   tag: stellar luminosity via Stefan-Boltzmann
Calculate the total luminosity of a star whose surface temperature is $7500$\,K,
and whose radius is $2.5$ times that of our Sun. Give your answer in units of the
solar luminosity, assuming the surface temperature of the Sun to be $5800$\,K.

\probhead{7.2}{}{2008}{Bandung, Indonesia}{TH S4}{20}
% source id: 2008_TH_S4   tag: helium-burning lifetime, triple-alpha
Suppose a star has a mass of $20\,M_\odot$. If 20\% of the star's mass is now in
the form of helium, calculate the helium-burning lifetime of this star. Assume
that the luminosity of the star is $100\,L_\odot$, of which 30\% is contributed
by helium burning. The carbon mass, $^{12}$C, is $12.000000$\,amu. Helium
burning to carbon: $3\,^{4}\mathrm{He} \rightarrow\ ^{12}\mathrm{C} + \gamma$.

\probhead{7.3}{}{2008}{Bandung, Indonesia}{TH S7}{20}
% source id: 2008_TH_S7   tag: giant-branch luminosity, visibility distance
A main sequence star at a distance $20$\,pc is barely visible through a certain
space-based telescope which can record all wavelengths. The star will eventually
move up along the giant branch, during which time its temperature drops by a
factor of 3 and its radius increases 100 times. What is the new maximum distance
at which the star can still be (barely) visible using the same telescope?

\probhead{7.4}{}{2010}{Beijing, China}{TH S12}{10}
% source id: 2010_TH_S12   tag: stellar radius ratio from dm and Teff
The difference in brightness between two main sequence stars in an open cluster
is 2 magnitudes. Their effective temperatures are $6000$\,K and $5000$\,K
respectively. Estimate the ratio of their radii.

\probhead{7.5}{}{2010}{Beijing, China}{TH S5}{10}
% source id: 2010_TH_S5   tag: maximum solar lifetime from mass-energy
If 0.8\% of the initial total solar mass could be transformed into energy during
the whole life of the Sun, estimate the maximum possible lifetime for the Sun.
Assume that the solar luminosity remains constant.

\probhead{7.6}{}{2012}{Rio de Janeiro, Brazil}{TH S7}{}
% source id: 2012_TH_S7   tag: LBV radius ratio at constant luminosity
Luminous Blue Variable (LBV) stars greatly vary in visual brightness; however,
the bolometric magnitude remains constant. Imagine an LBV star with a black body
temperature of $5000$\,K at its maximum visual brightness, and $30\,000$\,K at
its minimum visual brightness. Calculate the ratio of the star's radius between
the two situations above.

\probhead{7.7}{}{2013}{Volos, Greece}{TH S12}{10}
% source id: 2013_TH_S12   tag: Sirius B radius from flux ratio
Sirius A, with visual magnitude $m_V = -1.47$ (the brighter star on the sky) and
with stellar radius $R_A = 1.7\,R_\odot$, is the primary star of a binary system.
The existence of its companion, Sirius B, was deduced from astrometry in 1844 by
the well known mathematician and astronomer Friedrich Bessel, before it was
directly observed. Assuming that both stars were of the same spectral type and
that Sirius B is fainter by 10 magnitudes ($\Delta m = 10$), calculate the radius
of Sirius B.

\probhead{7.8}{}{2013}{Volos, Greece}{TH S3}{10}
% source id: 2013_TH_S3   tag: main-sequence lifetime of 5 Msun star
It is estimated that the Sun will have spent a total of about
$t_1 = 10$ billion years on the main sequence before evolving away from it.
Estimate the corresponding amount of time, $t_2$, if the Sun were 5 times more
massive.
\ioaafig{2013_TH_S3-f1.pdf}

\probhead{7.9}{}{2013}{Volos, Greece}{TH S6}{10}
% source id: 2013_TH_S6   tag: star distance, luminosity, radius
A star has an effective temperature $T_{\mathrm{eff}} = 8700$\,K, absolute
magnitude $M = 1.6$ and apparent magnitude $m = 7.2$. Find (a) the star's
distance $r$, (b) its luminosity $L$, and (c) its radius $R$. (Ignore
extinction.)

\probhead{7.10}{}{2013}{Volos, Greece}{TH S7}{10}
% source id: 2013_TH_S7   tag: red-giant classification via bolometric correction
A star has visual apparent magnitude $m_v = 12.2$\,mag, parallax
$\pi = 0''.001$ and effective temperature $T_{\mathrm{eff}} = 4000$\,K. Its
bolometric correction is $\mathrm{B.C.} = -0.6$\,mag.
(a) Find its luminosity as a function of the solar luminosity.
(b) What type of star is it: (i) a red giant, (ii) a blue giant, or (iii) a red
dwarf?

\probhead{7.11}{The density of the star}{2014}{Suceava, Romania}{TH P10}{}
% source id: 2014_TH_P10   tag: hydrostatic star, gas+radiation pressure
In a very simple model, a star is assumed to be a sphere of gas in a state of
equilibrium in its own gravitational field. The stellar gas consists of plasma,
i.e.\ hydrogen and helium atoms, completely ionized. Find an expression for the
value of the mass of the star if you know:
\begin{itemize}\itemsep1pt
\item $r$ --- the radius of the star;
\item $T$ --- the temperature of the star;
\item $n$ --- the relative proportion of hydrogen in the mass of the star;
\item $\mu_H$ --- the molar mass of hydrogen;
\item $\mu_{He}$ --- the molar mass of helium;
\item $R$ --- the universal gas constant;
\item $K$ --- the gravitational constant.
\end{itemize}
You may use the formula for the pressure of radiation inside the star,
\[ p_{rad} = \tfrac{1}{3}\,a T^4, \]
where $a$ is a known constant. The rotation of the star is negligible.

\probhead{7.12}{Life-time of a star on main sequence}{2014}{Suceava, Romania}{TH P6}{}
% source id: 2014_TH_P6   tag: MS lifetime from mass-luminosity
\ioaafig{2014_TH_P6-f1.pdf}
The plot of the function $\log(L/L_S) = f\!\left(\log(M/M_S)\right)$ for data
collected from a large number of stars is represented in the figure. The symbols
represent: $L$ and $M$ the luminosity and the mass of a star, and $L_S$ and $M_S$
the luminosity and the mass of the Sun respectively.

Find an expression for the lifetime of each star on the Main Sequence of the
Hertzsprung--Russell diagram, if the time spent by the Sun on the same Main
Sequence is $\tau_S$. Consider the following assumptions: for each star the
fraction of its mass converted into energy is $\eta$, the fraction of the mass
of the Sun converted into energy is $\eta_S$, the mass of each star is
$M = nM_S$, and the luminosity of each star remains constant during its entire
lifetime.

\probhead{7.13}{}{2015}{Magelang, Indonesia}{TH L3}{}
% source id: 2015_TH_L3   tag: star equation of state, heat capacities
Suppose a static spherical star consists of $N$ neutral particles with radius $R$
(see figure), with $0 \le \theta \le \pi$, $0 \le \phi \le 2\pi$, satisfying the
following equation of state:
\begin{equation}
  PV = N k \,\frac{T_R - T_0}{\ln(T_R/T_0)}
\end{equation}
where $P$ and $V$ are the pressure inside the star and the volume of the star
respectively, and $k$ is the Boltzmann constant. $T_R$ and $T_0$ are the
temperature at the surface $r = R$ and the temperature at the centre $r = 0$
respectively. Assume that $T_R \le T_0$.

\ioaafig{2015_TH_L3-f1.pdf}

\begin{description}
\item[(a)] Simplify the stellar equation of state (1) if
  $\Delta T = T_R - T_0 \approx 0$ (this is called an ideal star).
  (\textit{Hint:} use the approximation $\ln(1+x) \approx x$ for small $x$.)
\end{description}

Suppose the star undergoes a quasi-static process, in which it may slightly
contract or expand, such that the above stellar equation of state (1) still
holds. The star satisfies the first law of thermodynamics
\begin{equation}
  Q = \Delta M c^2 + W
\end{equation}
where $Q$, $M$ and $W$ are the heat, the mass of the star, and the work
respectively, while $c$ is the speed of light in vacuum and
$\Delta M \equiv M_{\mathrm{final}} - M_{\mathrm{initial}}$.

In the following we assume $T_0$ to be constant, while $T_R \equiv T$ varies.

\begin{description}
\item[(b)] Find the heat capacity of the star at constant volume $C_v$ in terms
  of $M$, and at constant pressure $C_p$ expressed in $C_v$ and $T$.
  (\textit{Hint:} use the approximation $(1+x)^n \approx 1 + nx$ for small $x$.)
\item[(c)] Assuming that $C_v$ is constant and the gas undergoes an isobaric
  process so the star produces heat and radiates it outside into space, find the
  heat produced by the isobaric process if the initial temperature and the final
  temperature are $T_i$ and $T_f$ respectively.
\item[(d)] Suppose there is an observer far away from the star. Using
  information from part (c), estimate the distance of the observer from the
  star.
\end{description}

For the next parts, assume the star is the Sun.

\begin{description}
\item[(e)] If the sunlight is monochromatic with frequency $5\times10^{14}$\,Hz,
  estimate the number of photons radiated by the Sun per second.
\item[(f)] Calculate the heat capacity $C_v$ of the Sun assuming its surface
  temperature varies from $5500$\,K to $6000$\,K in one second.
\end{description}

\probhead{7.14}{Expanding ring nebula}{2022}{Kutaisi, Georgia}{TH 3}{10}
% source id: 2022_TH_3   tag: planetary-nebula expansion, shell resolution
A planetary nebula, located $100$\,pc from the Earth, has the shape of a perfect
circular ring with an inner radius of $7.0'$ and an outer radius of $8.0'$. Its
luminosity is powered by the UV radiation from the white dwarf remnant at the
centre of the nebula. From other observations, we know that 2000 years ago the
inner radius of the nebula was $3.5'$ and the outer radius was $4.0'$. We believe
that throughout these 2000 years the evolution of the nebula obeys a free
expansion scenario, so gravity is negligible and the expansion velocity remains
constant with time. Assume that all material in the planetary nebula was ejected
at the same instant, but different gas particles have different velocities.

\begin{description}
\item[(a)] Estimate the range of velocities of the gas particles. \hfill[4\,pt]
\item[(b)] Is the assumption of free expansion justified? Write YES or NO
  alongside appropriate calculations. \hfill[3\,pt]
\item[(c)] If this planetary nebula is bright enough, would an astronaut aboard
  the International Space Station be able to resolve the thickness of the shell
  clearly? Write YES or NO alongside the necessary calculations. \hfill[3\,pt]
\end{description}

\probhead{7.15}{Flaring protoplanetary disk}{2022}{Kutaisi, Georgia}{TH 5}{10}
% source id: 2022_TH_5   tag: flaring disk thermal structure
Planets are the products of collapsing clouds forming circumstellar disks around
young stars. In this problem, we examine the thermal structure of a type of
protoplanetary disk. We consider radiation from the central star to be the
dominant heating process and assume that the disk is optically thick and the
radiation can only be absorbed in the surface layers of the disk. This type of
disk is called a flaring disk (see figure below). In this case, the star
illuminates the surface of the disk directly, and the shallow angle $\beta$
between the light rays and the surface of the disk increases with distance from
the star.

\begin{description}
\item[(a)] Find an expression for $\beta$ as a function of $r$. Assume
  $h(r) \ll r$. \hfill[4\,pt]\\
  \textit{Hint:} the expression involves $h(r)$, $r$ and $\Delta h(r)/\Delta r$.
\item[(b)] Find an expression for the equilibrium temperature of the disk
  ($T_D$) as a function of $\beta(r)$, $r$ and the luminosity ($L_s$) of the star
  (ignore the size of the star). \hfill[3\,pt]
\item[(c)] We parameterize the disc height as $h(r) = a r^b$, where $a$ and $b$
  are constants. For which values of $a$ and $b$ is the condition for an
  isothermal layer ($T_D = $ constant) satisfied? Write your answer for the
  constant $a$ in terms of $T_D$ and $L_S$. \hfill[3\,pt]\\
  \textit{Hint:} if $\Delta r \ll r$, then
  $(1 + \Delta r/r)^b \approx 1 + b\,\Delta r/r$.
\end{description}
\ioaafig{2022_TH_5-f1.pdf}

\probhead{7.16}{Neutrinos}{2023}{Katowice, Poland}{TH Q8}{20}
% source id: 2023_TH_Q8   tag: supernova vs solar neutrino flux
In a simplified model of a supernova explosion, the core of a star, composed of
pure iron $^{56}_{26}$Fe nuclei with a total mass of $1\,M_\odot$, changes into a
neutron star composed of individual electrons, protons and neutrons in numerical
proportions of $1:1:8$. This process is called ``neutronization'' and results in
the emission of a large number of neutrinos.

Calculate the solar neutrino flux on Earth. How much larger would the flux of
neutrinos reaching the Earth from the supernova be than the steady neutrino
emission of the Sun, if the supernova exploded in the centre of the Galaxy and
the process of neutronization of the core took about $0.01$\,s? Give an
order-of-magnitude answer.

\probhead{7.17}{Stars through graphs}{2025}{Mumbai, India}{TH T11}{50}
% source id: 2025_TH_T11   tag: stellar interior profiles, evolution
\ioaafig{2025_TH_T11-f1.pdf}
Stars can be well approximated as spherically symmetric objects, and hence the
radial distance $r$ from the centre can be chosen as the only independent
variable in modelling stellar interiors. The mass contained within a sphere of
radius $r$ is denoted by $m(r)$. The luminosity $l(r)$ is defined as the net
energy flowing outward through a spherical surface of radius $r$ per unit time.
Other quantities of interest, for example, the density $\rho(r)$, temperature
$T(r)$, hydrogen mass fraction $X(r)$, helium mass fraction $Y(r)$, and the
nuclear energy generated per unit mass per unit time
$\epsilon_{\mathrm{nuc}}(r)$, are taken to be functions of $r$. Throughout this
problem we shall neglect the effects of diffusion and gravitational settling of
elements inside the star.

The symbol ``log'' refers to logarithm to the base 10. The problem consists of
three independent parts.

\medskip\noindent\textbf{Part 1: Inside a star}

The graph below shows the variation of three structural quantities, A, B, and
C, as functions of the fractional radius $r/R$ in a stellar model of mass
$1\,\mathrm{M_\odot}$ and age 4\,Gyr, where $R$ is the photospheric radius of
the star. The values of the helium mass fraction at the (photospheric) surface,
$Y_{\mathrm{s}}$, and the metallicity (mass fraction of all elements heavier
than helium) at the (photospheric) surface, $Z_{\mathrm{s}}$, of the star are
given by $(Y_{\mathrm{s}},\,Z_{\mathrm{s}}) = (0.28,\,0.02)$. All quantities
shown in the plots are normalised by their respective maximum values.
% MISSING FIGURE: normalized quantities A, B, C vs r/R (2025 paper p.13) -- needs extraction into images/ (e.g. 2025_TH_T11-f0.pdf)

\begin{description}
\item[(T11.1a)] Identify the three quantities A, B, and C uniquely from among
  the five possibilities:
  \[ T(r),\ l(r),\ \epsilon_{\mathrm{nuc}}(r),\ X(r),\ Y(r). \]
  (Write A/B/C in the boxes beside the appropriate quantities in the Summary
  Answersheet. No justification is needed for your answer.) \hfill[6\,pt]
\item[(T11.1b)] What is the mass fraction of helium at the centre,
  $Y_{\mathrm{c}}$, of the star? \hfill[3\,pt]
\item[(T11.1c)] Sketch the remaining two quantities from the list of five
  (which were not identified as curves A, B, or C) given in (T11.1a), as
  functions of $r/R$ on the same graph in the Summary Answersheet, and label by
  their respective quantities. \hfill[5\,pt]
\end{description}

\medskip\noindent\textbf{Part 2: Evolving stars}

Consider the evolution of a $1\,\mathrm{M_\odot}$ star whose initial uniform
composition is given by the mass fractions of helium, $Y_0 = 0.28$, and metals,
$Z_0 = 0.02$. The figures below show the variation of different global
quantities of this star as it evolves from ZAMS (Zero Age Main Sequence) till
the end of helium burning in its core.

The graph below shows the evolutionary track of the star on the HR diagram
(plot of $\log L/L_\odot$ vs $\log T_{\mathrm{eff}}$, where $L$ is the surface
luminosity and $T_{\mathrm{eff}}$ is the effective temperature).
\ioaafig{2025_TH_T11-f1.pdf}

The figure below has four graphs which show the variation of
$T_{\mathrm{eff}}$ (in K), $L$ (plotted as $\log L/L_\odot$), $R$ (plotted as
$\log R/R_\odot$), and $Y_{\mathrm{c}}$ with age (in $10^9$\,yr) of the same
star. In each of these four graphs, the insets show the variations of the
respective quantities in detail between the ages of $11.86\times10^9$\,yr to
$12.00\times10^9$\,yr, for greater clarity.
\ioaafig{2025_TH_T11-f2.pdf}
\ioaafig{2025_TH_T11-f3.pdf}

Use these graphs to answer the questions below.
\begin{description}
\item[(T11.2a)] What is the approximate main sequence lifetime,
  $t_{\mathrm{MS}}$ (in years), of the star? \hfill[1\,pt]
\item[(T11.2b)] What is the approximate duration, $\Delta t_{\mathrm{He}}$ (in
  years), for which the star burns helium in its core? \hfill[1\,pt]
\item[(T11.2c)] What fraction, $f_{\mathrm{H}}$, of the initial amount of
  hydrogen at its centre has been burnt when the luminosity of the star is
  $1\,\mathrm{L_\odot}$? \hfill[3\,pt]
\item[(T11.2d)] What is the radius of the star, $R_1$ (in units of
  $\mathrm{R_\odot}$) when 60\% of the initial amount of hydrogen at its centre
  has been burnt? \hfill[3\,pt]
\item[(T11.2e)] What are the radii of the star, $R_{\mathrm{P}}$ and
  $R_{\mathrm{Q}}$ (in units of $\mathrm{R_\odot}$), corresponding to its
  positions P and Q, respectively, as marked on the HR-diagram? \hfill[4\,pt]
\end{description}

\medskip\noindent\textbf{Part 3: Mass distribution inside a star}

The equation that governs the distribution of mass inside a star is given by
\[ \frac{dm(r)}{dr} = 4\pi r^2 \rho(r) \]
It would be convenient to express this equation in terms of three dimensionless
variables, namely, the fractional mass, $q$, the fractional radius, $x$, and
the relative density, $\sigma$, that we define as
\[ q = m/M \qquad x = r/R \qquad \sigma = \rho/\bar{\rho} \]
where $M$ and $R$ are the total mass and radius of the star, respectively, and
$\bar{\rho} \equiv \dfrac{M}{\frac{4}{3}\pi R^3}$ is the average density of the
star. For the particular star that we shall be considering in this part, the
following information is given:
\begin{itemize}
\item The central density $\rho(x = 0) = 80\bar{\rho}$
\item Half of the star's mass is contained within the inner 25\% of its total
  radius, and 70\% of its mass is contained within the inner 35\% of its total
  radius.
\end{itemize}
In all subsequent parts of this question, it will be sufficient to round off
all derived numerical coefficients to within 0.005.
\begin{description}
\item[(T11.3a)] Express the above equation describing the dependence of mass on
  radius in terms of $x$, $\dfrac{dq(x)}{dx}$ and $\sigma(x)$. \hfill[2\,pt]
\end{description}

To obtain the distribution of mass with radius, we need to know the density
profile inside the star. For the purpose of this problem, we shall describe the
variation of density with radius by approximate forms in two domains of $x$:
\begin{itemize}
\item the inner part of the star: $0 \le x \le 0.32$
\item the middle part of the star: $0.32 < x < 0.80$
\end{itemize}
We do not make any approximation for the outermost part, i.e.,
$0.80 \le x \le 1.00$.
\begin{description}
\item[(T11.3b)] \textbf{Approximation for the middle part}: The variation of
  $\log\sigma$, as a function of $\log x$ in the middle part of the star is
  shown (by the black curve) in the graph below. We shall make a linear
  approximation (shown as a dashed red line in the graph) for $\log\sigma$ as a
  function of $\log x$ in the domain $-0.5 < \log x < -0.1$, i.e.,
  $0.32 \lesssim x \lesssim 0.80$ (shown by the green shaded domain). Further,
  we shall approximate the slope of this line by the nearest integer.
  \ioaafig{2025_TH_T11-f4.pdf}
  Use this approximation to write an expression for $\sigma(x)$ as a function
  of $x$ in the domain $0.32 < x < 0.80$. \hfill[4\,pt]
\item[(T11.3c)] Use the result of (T11.3b) to derive an expression for $q(x)$
  in the domain $0.32 < x < 0.80$. \hfill[6\,pt]
\item[(T11.3d)] \textbf{Approximation for the inner part}: In the inner part of
  the star ($0 \le x \le 0.32$), the density may be approximated as a linear
  function of radius, i.e., $\sigma(x) = Ax + B$, where $A, B$ are constants.
  Determine $A$ and $B$, and hence obtain an expression for $q(x)$ in the
  domain $0 \le x \le 0.32$. Note that the approximations adopted in the
  previous part and this part may lead to small discontinuities in density or
  mass at $x = 0.32$. \hfill[8\,pt]
\item[(T11.3e)] The expressions for $q(x)$ obtained in parts (T11.3c) and
  (T11.3d) are approximations that describe the variation of mass with radius
  quite well, but only in specific regions of the star. For the domain
  $0.80 \le x \le 1$ (for which we have not derived any expression), it is
  possible to use appropriate extrapolation from the neighbouring region. Use
  these approximate expressions and given data to sketch a smooth curve
  (without any discontinuities either in $q(x)$ or its derivative) for $q(x)$
  vs $x$ for the entire star ($0 \le x \le 1$) that represents the variation of
  mass with radius. \hfill[4\,pt]
\end{description}

\clearpage
\section*{Solutions to Chapter 7}

\solhead{7.1}{}
% source id: 2007_TH_1.9
From the Stefan--Boltzmann law, the total power emitted is
\[
  L = \left(4\pi R^2\right)\left(\sigma T^4\right)
    = \left(4\pi R_\odot^2\right)\left(\sigma T_\odot^4\right)
      \left(\frac{R}{R_\odot}\right)^{\!2}\left(\frac{T}{T_\odot}\right)^{\!4}
    = L_\odot \left(2.5\right)^2 \left(\frac{7500}{5800}\right)^{\!4}
\]
\begin{align*}
  \frac{L}{L_\odot} &= \left(6.25\right)\left(2.796\right) = 17.47 = 17.5 \\
  L &= 17.5\,L_\odot
\end{align*}

\solhead{7.2}{}
% source id: 2008_TH_S4
The mass of three helium nuclei $= 3 \times 4.002603$ amu $= 12.007809$ amu
\hfill(20\%)

The mass converted to energy is the difference between the sum of the masses of
three helium nuclei and the mass of the resulting carbon nucleus:
\[
  12.007809\ \text{amu} - 12.000000\ \text{amu} = 7.81 \times 10^{-3}\ \text{amu}.
\]
\hfill(15\%)

This represents a fractional loss of mass of
$\dfrac{7.81 \times 10^{-3}}{12.007809} = 6.5 \times 10^{-4}$ of the original
mass.
\hfill(15\%)

The stellar mass $m = 20 \times \left(1.99 \times 10^{33}\right) = 3.98 \times 10^{34}\,\mathrm{g}$

Thus $6.5 \times 10^{-4}$ of 20\% of the star's mass will be converted into
energy during the triple-$\alpha$ process, or
\[
  E = mc^2 = 0.2(3.98 \times 10^{34})(6.5 \times 10^{-4})(3 \times 10^{10})^2
    = 4.66 \times 10^{51}\,\mathrm{erg}
\]
\hfill(25\%)

The helium-burning lifetime is
\[
  t = \frac{E}{0.3L}
    = \frac{4.66 \times 10^{51}}{0.3(100)(3.96 \times 10^{33})}
    = 3.92 \times 10^{16}\,\mathrm{s} = 1.24 \times 10^{9}\ \text{years}
\]
\hfill(25\%)

\solhead{7.3}{}
% source id: 2008_TH_S7
Let $L$ = luminosity of the giant stars in the giant branch,
$L_{MS}$ = luminosity of the main sequence star,
$F$ = flux of the giant stars in the giant branch,
$F_{MS}$ = flux of the main sequence star,
$T$ = temperature of the giant star,
$T_{MS}$ = temperature of the main sequence star,
$R$ = radius of the giant star,
$R_{MS}$ = radius of the main sequence star.

Use the radius--luminosity--temperature relation
\[
  \frac{L}{L_{MS}} = \frac{R^2 T^4}{R_{MS}^2 T_{MS}^4}
    = \left(\frac{R}{R_{MS}}\right)^{\!2}\left(\frac{T}{T_{MS}}\right)^{\!4}
    = \left(100\right)^2 \left(\frac{1}{3}\right)^{\!4} = 123.5
\]
\hfill(50\%)
\begin{align*}
  F &= \frac{L}{4\pi d^2} \\
  \frac{F}{F_{MS}} &= \frac{L}{L_{MS}}\left(\frac{d_{MS}}{d}\right)^{\!2}
\end{align*}
that is, the star becomes 123.5 times brighter when it becomes a giant star. To
be just barely visible using the same telescope, we push the star farther so
that it becomes 123.5 times dimmer (flux is inversely proportional to a square
of the distance).
\begin{align*}
  \frac{1}{123.5} &= \left[\frac{1/d}{1/d_{MS}}\right]^{2}
  \;\Longrightarrow\; \frac{d}{d_{MS}} = 11.11
  \;\Longrightarrow\; d = 11.11\,d_{MS} \\
  &= 11.11\,(20)\ \text{pc} = 222.20\ \text{pc}
\end{align*}
\hfill(50\%)

\solhead{7.4}{}
% source id: 2010_TH_S12
\[
  L_1 = 4\pi R_1^2 \sigma T_{\max}^4
\]
\hfill(3 points)
\[
  L_2 = 4\pi R_2^2 \sigma T_{\min}^4
\]
\[
  \Delta m = -2.5\lg(L_{\min}/L_{\max})
           = -5\lg(R_{\min}/R_{\max}) - 10\lg(T_{\min}/T_{\max})
\]
\hfill(3 points)
\[
  \lg(R_{\min}/R_{\max}) = -0.2\,\Delta m - 2\lg(T_{\min}/T_{\max}) = -0.24
\]
\hfill(2 points)

So,
\[
  R_{\min}/R_{\max} = 0.57
\]
\hfill(2 points)

\solhead{7.5}{}
% source id: 2010_TH_S5
The total mass of the Sun is
\[
  m \approx 1.99 \times 10^{30}\,\mathrm{kg}
\]
0.8\% mass transform into energy:
\[
  E = mc^2 \approx 0.008 \times 2 \times 10^{30} \times (3 \times 10^{8})^2
    = 1.4 \times 10^{45}\,\mathrm{J}
\]
\hfill(5 points)

Luminosity of the Sun is
\[
  L_{\mathrm{sun}} = 3.96 \times 10^{26}\,\mathrm{W}
\]
Sun's life would at most be:
\[
  t = E / L_{\mathrm{sun}} = 3.6 \times 10^{18}\,\mathrm{s} \approx 10^{11}\ \text{years}
\]
\hfill(5 points)

\solhead{7.6}{}
% source id: 2012_TH_S7
\[
  T_1 = 30000\ \mathrm{K}
\]
\hfill(4 pt)
\[
  T_2 = 5000\ \mathrm{K}
\]
\[
  m_{\mathrm{bol}} = \text{constant} \;\Longleftrightarrow\; L = \text{constant}
\]
\hfill(4 pt)

As $L = 4\pi\sigma R^2 T^4$ and $L_1 = L_2$, we have:
\[
  4\pi\sigma R_0^2 T_1^4 = 4\pi\sigma R_{\mathrm{shell}}^2 T_2^4
  \;\Longrightarrow\; \frac{R_{\mathrm{shell}}}{R_0} = 36
\]
\hfill(2 pt)

\solhead{7.7}{}
% source id: 2013_TH_S12
The distance of the two stars from our solar system is the same. Therefore
\[
  m_B - m_A = 2.5\,\log\frac{\dfrac{L_A}{4\pi r^2}}{\dfrac{L_B}{4\pi r^2}}
            = 2.5\,\log\frac{L_A}{L_B}
\]
\hfill(5 Point)

From which we get $L_A = 10^4\,L_B$. From equation
$L = 4\pi R^2 \sigma T_{\mathrm{eff}}^4$ we get
\[
  L_A/L_B = (R_A/R_B)^2\,(T_A/T_B)^4
\]
Assuming that the two stars belong to the same spectral type (and therefore
$T_A = T_B$) we get
\[
  R_B = 0.01 R_A = 0.01 \times 1.7 \times 696000\ \mathrm{km}
  \quad\text{or}\quad R_B = 1.2\times10^{4}\ \mathrm{km}
\]
\hfill(5 Point)

\solhead{7.8}{}
% source id: 2013_TH_S3
For the average luminosity of a main sequence star we have: $L \propto M^4$
(where $M$ is the initial mass of the star). We assume that the total energy
$E$ that the star produces is proportional to its mass, $E \propto M$.
Therefore the amount of time that the star spends on the main sequence is
approximately
\[
  t_{MS} \approx \frac{E}{L} \propto \frac{M}{M^4} \approx M^{-3}.
\]
\hfill(5 Points)

Therefore,
\[
  \frac{t_1}{t_2} \approx \frac{1^{-3}}{5^{-3}} = 5^3
  \;\rightarrow\;
  t_2 = \frac{1}{125}\times10^{10}\ \mathrm{yr}
  \quad\text{or}\quad t_2 = 8\times10^{7}\ \mathrm{yr}.
\]
\hfill(5 Points)

\solhead{7.9}{}
% source id: 2013_TH_S6
\begin{description}
\item[(a)] Its distance is calculated from equation: $m - M = 5\,\log(r) - 5$,
  or $7.2 - 1.6 + 5 = 5\log(r) \;\rightarrow\; \log(r) = 2.12$ and
  $r = 132\ \mathrm{pc}$
  \hfill(3 Points)
\item[(b)] Its Luminosity is calculated from equation:
  $M_\odot - M = 2.5\,\log\!\left(\dfrac{L}{L_\odot}\right)$ or
  \[
    4.8 - 1.6 = 2.5\,\log\!\left(\frac{L}{L_\odot}\right)
    \ \text{or}\
    \log\!\left(\frac{L}{L_\odot}\right) = 1.28
    \;\rightarrow\;
    \left(\frac{L}{L_\odot}\right) = 19.05
    \ \text{or}\
    L = 19.15\times L_\odot % sic: 19.15 vs 19.05 in source
    = 19.15\times3.9\times10^{33}
  \]
  and $L = 7.4\times10^{34}\ \mathrm{erg\,sec^{-1}}$
  \hfill(4 Points)
\item[(c)] Its radius can be easily calculated from equation:
  $L = 4\pi\sigma R^2\,T_{\mathrm{eff}}^4$, from which
  \[
    R = \frac{1}{T_{\mathrm{eff}}^2}\sqrt{\frac{L}{4\pi\sigma}}
  \]
  from which we get: $R = 1.35\times10^{11}\ \mathrm{cm}$
  \hfill(3 Points)
\end{description}

\solhead{7.10}{}
% source id: 2013_TH_S7
\begin{description}
\item[(a)] First its bolometric magnitude is calculated from equation:
  $M_V - m_V = 5 - 5\,\log(r)$ or equivalent:
  $M_V - m_V = 5 + 5\,\log\pi \;\rightarrow\;
  M_V = 12.2 + 5 + 5\log(0''.001) = 12.2 + 5 - 15 = 2.2$~mag.
  Its barycentric correction is: % sic: source says "barycentric" for bolometric
  $B.C. = M_{\mathrm{bol}} - M_V$ and $M_{\mathrm{bol}} = B.C. + M_V$ or
  $M_{\mathrm{bol}} = -0.6 + 2.2$ or $M_{\mathrm{bol}} = 1.6$~mag.
  \hfill(4 Points)

  Then its Luminosity is calculated from:
  \[
    M_\odot - M_{\mathrm{bol}} = 2.5\,\log\!\left(\frac{L}{L_\odot}\right),
    \ \text{or}\
    4.72 - 1.6 = 2.5\,\log\!\left(\frac{L}{L_\odot}\right)
    \ \text{or}\
    \log\!\left(\frac{L}{L_\odot}\right) = 1.25
    \ \text{and}\ L = 17.70\,L_\odot
  \]
  \hfill(2 Point)
\item[(b)] Type of star: A star with $M_{\mathrm{bol}} = 1.6$~mag,
  $L = 17.7\,L_\odot$ and $T_{\mathrm{eff}} = 4000$~K is much brighter and much
  cooler than the Sun (see \emph{Table of constants}). Therefore it is (i) a red
  giant star.
  \hfill(4 Points)
\end{description}

\solhead{7.11}{The density of the star}
% source id: 2014_TH_P10
The hydrostatic equilibrium inside the star means that in each point of the
inside of the star the gravitational forces are compensated by the hydrostatic
pressure forces. That means that the mater of the star remains localized in a
region of space.

The total pressure of the stellar gas has two components: the pressure due to
the movement of the stellar-gas particles $\left(p_{\mathrm{gaz}}\right)$ and
the pressure due to the emitted radiation by the stellar-gas particles
$\left(p_{\mathrm{rad}}\right)$, thus:
\begin{align*}
  p_{\mathrm{total}} &= p_{\mathrm{gaz}} + p_{\mathrm{rad}}; \\
  p_{\mathrm{gaz}} &= p_{\mathrm{H}} + p_{\mathrm{He}}; \\
  p_{\mathrm{rad}} &= \frac{1}{3}aT; % sic: source writes aT here and below (should be aT^4)
\end{align*}
\[
  p_{\mathrm{total}}
  = \rho\,\frac{n\mu_{\mathrm{He}} + (1-n)\mu_{\mathrm{H}}}{\mu_{\mathrm{H}}\mu_{\mathrm{He}}}\,RT
    + \frac{1}{3}aT. \tag{1}
\]
In order to calculate the gravitational pressure of the stellar gas let's
consider a narrow cylinder with section area $\Delta S$, along the radius of
the star. See the figure below. If the total gravitational force acting on this
cylinder is $\vec{F}_g$ than the gravitational pressure exert by the
gas-column is $p_{\mathrm{grav}} = F_g/\Delta S$.
\ioaafig{2014_TH_P10-s1.pdf}

In order to calculate $\vec{F}_g$ let divide the cylinder in $n$ identically
small cylinders each of it with height $\Delta r$ and mass $\Delta m$.
Considering an homogenous star:
\begin{align*}
  F_1 &= K\frac{\Delta m \cdot M_1}{\left(r - \dfrac{\Delta r}{2}\right)^{2}}
       = K\frac{\Delta m \cdot M_1}{r^2\left(1 - \dfrac{\Delta r}{2r}\right)^{2}}
       = K\frac{\Delta m \cdot M_1}{r^2}\left(1 - \frac{\Delta r}{2r}\right)^{\!-2}; \\
  \frac{M}{r^3} &= \frac{M_1}{(r-\Delta r)^3}; \quad
  M_1 = M\cdot\frac{(r-\Delta r)^3}{r^3} = M\cdot\left(1 - \frac{\Delta r}{r}\right)^{\!3}; \\
  F_1 &\approx K\frac{\Delta m \cdot M}{r^2}\left(1 - 2\frac{\Delta r}{r}\right); \\
  F_2 &\approx K\frac{\Delta m \cdot M}{r^2}\left(1 - 3\frac{\Delta r}{r}\right); \\
  &\dotso \\
  F_n &\approx K\frac{\Delta m \cdot M}{r^2}\left(1 - (n+1)\frac{\Delta r}{r}\right); \\
  F_g &= F_1 + F_2 + \dotsb + F_n; \\
  F_g &= K\frac{m\cdot M}{n r^2}\left[n - \frac{1}{n}\,\frac{(1+n)n}{2}\right]; \quad
  F_g = K\frac{m\cdot M}{r^2}\left[1 - \frac{(1+n)}{2n}\right]; \\
  F_g &= K\frac{m\cdot M}{r^2}\,\frac{n-1}{2n}; \quad
  \frac{n-1}{2n} \approx \frac{1}{2}; \\
  F_g &= K\frac{m\cdot M}{2r^2}. \tag{2}
\end{align*}
Thus the gravitational pressure is
\begin{align*}
  p_g &= \frac{F_g}{\Delta S}
       = \frac{1}{\Delta S}\,
         \frac{\rho\cdot r\cdot\Delta S\cdot\rho\cdot\dfrac{4\pi r^3}{3}}{2r^2} \\
  p_g &= \frac{2}{3}\pi K r^2 \rho^2
\end{align*}
Using the relations (1) and (2) in the pressures equilibrium relationship
\[
  p_{\mathrm{total}} = p_{\mathrm{gravitational}}
\]
Results:
\begin{align*}
  \frac{2}{3}\pi K r^2\rho^2
    &= \rho\,\frac{n\mu_{\mathrm{He}} + (1-n)\mu_{\mathrm{H}}}{\mu_{\mathrm{H}}\mu_{\mathrm{He}}}\,RT
       + \frac{1}{3}aT; \\
  \frac{2}{3}\pi K r^2\cdot\rho^2
    - \frac{n\mu_{\mathrm{He}} + (1-n)\mu_{\mathrm{H}}}{\mu_{\mathrm{H}}\mu_{\mathrm{He}}}\,RT\cdot\rho
    - \frac{1}{3}aT &= 0;
\end{align*}
This is an second degree equation in $\rho$
\[
  A\rho^2 - B\rho - C = 0
\]
Where the coefficients are
\begin{align*}
  A &= \frac{2}{3}\pi K r^2 \\
  B &= \frac{n\mu_{\mathrm{He}} + (1-n)\mu_{\mathrm{H}}}{\mu_{\mathrm{H}}\mu_{\mathrm{He}}}\,RT \\
  C &= \frac{1}{3}aT
\end{align*}
The positive solution is the valid one i.e.
\[
  \rho = \frac{B + \sqrt{B^2 + 4AC}}{2A}, \qquad
  \rho = \frac{\dfrac{n\mu_{\mathrm{He}} + (1-n)\mu_{\mathrm{H}}}{\mu_{\mathrm{H}}\mu_{\mathrm{He}}}\,RT
    + \sqrt{\left(\dfrac{n\mu_{\mathrm{He}} + (1-n)\mu_{\mathrm{H}}}{\mu_{\mathrm{H}}\mu_{\mathrm{He}}}\,RT\right)^{\!2}
    + \dfrac{8}{9}\pi K r^2 aT}}
    {\dfrac{4}{3}\pi K r^2}
\]

\solhead{7.12}{Life-time of a star on main sequence}
% source id: 2014_TH_P6
\textbf{A. The analysis of the graph:}

The graph is linear:
\[
  y = ax + b = ax;
\]
From the graph it can be obtain the following data:
\begin{align*}
  \log\frac{L}{L_S} &= a\cdot\log\frac{M}{M_S}; \\
  a = \tan\alpha &= \frac{\Delta y}{\Delta x} = \frac{3.5}{1} = 3.5;
\end{align*}
\ioaafig{2014_TH_P6-s1.pdf}
\[
  L \sim M^{3.5}. \tag{1}
\]
The total energy of the star is:
\[
  E = Mc^2,
\]
So the emitted energy due to the mass variation of the star is:
\[
  \Delta E = c^2 \Delta M,
\]
According to the text
\begin{align*}
  \Delta M &= \eta M; \\
  \Delta E &= c^2 \eta M.
\end{align*}
By using the definition of the luminosity:
\begin{align*}
  \frac{\Delta E}{\Delta t} &= L; \tag{2} \\
  \Delta t &= \tau; \\
  \frac{c^2\eta M}{\tau} &= L; \\
  \tau &= \frac{c^2\eta M}{L}, \tag{2}
\end{align*}
Which represents the life-time of the star.

By using the results from the graph analysis
\[
  L = \frac{L_S}{M_S^{3.5}}\cdot M^{3.5},
\]
Thus:
\[
  \tau = \frac{c^2\eta M_S^{3.5}}{L_S}\cdot M^{-2.5}.
\]
If use the same calculations for the Sun it can be obtain
\[
  E_S = M_S c^2;
\]
\[
  \tau_S = \frac{c^2\eta_S M_S}{L_S},
\]
Which is the life-time of the Sun
\begin{align*}
  \tau &= \frac{\eta}{\eta_S}\cdot\tau_S\cdot M_S^{2.5}\cdot M^{-2.5}; \\
  \frac{\tau}{\tau_S} &= \frac{\eta}{\eta_S}\cdot\left(\frac{M}{M_S}\right)^{\!-2.5};
  \quad M = nM_S; \\
  \tau &= \frac{\eta}{\eta_S}\,(n)^{-2.5}\,\tau_S.
\end{align*}

\solhead{7.13}{}
% source id: 2015_TH_L3
\begin{description}
\item[a.] Defining $\Delta T = T_R - T_0$ and $\Delta T \approx 0$, we have
  \[
    PV = Nk\,\frac{\Delta T}{\ln(1 + \Delta T/T_0)} \tag{3}
  \]
  using $\ln(1 + \Delta T/T_0) \approx \Delta T/T_0$, we then obtain
  \[
    PV = NkT_0 \tag{4}
  \]
  \hfill(15 points)
\item[b.] The internal energy of the star is $U = Mc^2$ ($U(T) = M(T)c^2$ for
  ideal star). Thus, the constant volume heat capacity of the star has the form
  \[
    C_v = \left(\frac{\Delta Q}{\Delta T}\right)_{\!V}
        = \left(\frac{\Delta M}{\Delta T}\right)_{\!V} c^2
  \]
  for small $\Delta T$. Then, using first law of themodynamics, the constant
  pressure heat capacity of the star is
  \[
    C_p = \left(\frac{\Delta Q}{\Delta T}\right)_{\!P}
        = \left(\frac{\Delta M}{\Delta T}\right)_{\!V} c^2 + P\,\frac{\Delta V}{\Delta T}
        = C_v + P\,\frac{\Delta V}{\Delta T}
  \]
  for small $\Delta T$. Defining $\Delta T = T_2 - T_1$, then
  \[
    P\,\Delta V = Nk\left(
      \frac{T_1 - T_0 + \Delta T}{\ln\!\big((T_1 + \Delta T)/T_0\big)}
      - \frac{T_1 - T_0}{\ln(T_1/T_0)}\right)
  \]
  Using the approximation
  \begin{align*}
    \ln\!\big((T_1 + \Delta T)/T_0\big)
      &\approx \ln\!\left(\frac{T_1}{T_0}\right) + \frac{\Delta T}{T_1} \\
    \left(1 + \frac{\Delta T}{T_1\ln(T_1/T_0)}\right)^{\!-1}
      &\approx 1 - \frac{\Delta T}{T_1\ln(T_1/T_0)}
  \end{align*}
  then we have
  \[
    P\,\frac{\Delta V}{\Delta T}
      = \frac{Nk}{\ln(T/T_0)}\left(1 - \frac{(T - T_0)/T}{\ln(T/T_0)}\right)
  \]
  where $T_1 \equiv T$. Finally, we obtain
  \[
    C_p = \left(\frac{\Delta Q}{\Delta T}\right)_{\!P}
        = \left(\frac{\Delta M}{\Delta T}\right)_{\!V} c^2 + P\,\frac{\Delta V}{\Delta T}
        = C_v + \frac{Nk}{\ln(T/T_0)}\left(1 - \frac{(T - T_0)/T}{\ln(T/T_0)}\right)
  \]
  \hfill(30 points)
\item[c.] Since $C_V$ is constant, the heat produced by the star is given by
  \begin{align*}
    Q_H &= C_v\left(T_f - T_i\right) + P\,\Delta V \\
        &= C_v\left(T_i - T_f\right) % sic: source flips to (T_i - T_f) in the second line
         + Nk\left(\frac{T_f - T_0}{\ln\!\big(T_f/T_0\big)}
         - \frac{T_i - T_0}{\ln(T_i/T_0)}\right)
  \end{align*}
  \hfill(20 points)
% Part d. is blank in the official solutions (15 points allotted), although the problem paper does pose a part d.
\item[e.] Energy per second radiated by the Sun
  \[
    L_\odot = Nh\nu
  \]
  where $N$ is the number of photon. Thus
  \[
    N = \frac{L_\odot}{h\nu}
      = \frac{3.96\times10^{26}}{6.6261\times10^{-34}\times 5\times10^{14}}
      = 1.195\times10^{45}\ \text{photons}
  \]
  \hfill(10 points)
\item[f.] Energy per second radiated by the Sun is proportional to mass defect
  of the Sun
  \[
    L_\odot = \Delta M\,c^2
  \]
  Thus,
  \[
    C_v = \frac{\Delta M\,c^2}{\Delta T} = \frac{L_\odot}{\Delta T}
        = \frac{3.96\times10^{26}}{6000 - 5500}\ \mathrm{J/K}
        = 7.92\times10^{23}\ \mathrm{J/K}
  \]
  \hfill(10 points)
\end{description}

\solhead{7.14}{Expanding ring nebula}
% source id: 2022_TH_3
\begin{description}
\item[(a)] The distance traveled by the gas particles of inner and outer radius
  can be calculated as follows:
  \begin{align*}
    S_{\mathrm{out}} &= d_\oplus(\theta'_{\mathrm{out}} - \theta_{\mathrm{out}})
      = 100\,\mathrm{pc}\cdot(8.0' - 4.0') = 24\times10^{3}\ \mathrm{au} \\
    S_{\mathrm{in}} &= d_\oplus(\theta'_{\mathrm{in}} - \theta_{\mathrm{in}})
      = 100\,\mathrm{pc}\cdot(7.0' - 3.5') = 21\times10^{3}\ \mathrm{au}
  \end{align*}
  \hfill(2.0 points)
  \begin{align*}
    \therefore\; v_{max} &= \frac{S_{\mathrm{out}}}{t_0} = 57.0\ \mathrm{km/s} \\
    v_{min} &= \frac{S_{\mathrm{in}}}{t_0} = 50.0\ \mathrm{km/s}
  \end{align*}
  \hfill(2.0 points)
\item[(b)] At the inner edge, escape velocity from the white dwarf can be
  maximal ($1.44M_\odot$ -- Chandrasekhar limit)
  \[
    v_{esc} = \sqrt{\frac{2G\cdot1.44M_\odot}{S_{\mathrm{in}}}}
      \approx 0.35\ \mathrm{km/s}
  \]
  \hfill(2.0 points)

  Since $v_{esc} \ll v_{min}$, the no-gravity scenario can be applied. YES
  \hfill(1.0 points)
\item[(c)]
  \begin{align*}
    \phi &= 8' - 7.5' = 0.5' = 30'' \ge \frac{1.22\lambda}{D} \\
    \therefore\; D &\ge \frac{1.22\lambda}{\phi}
      = \frac{1.22\times5\times10^{-7}\times206265}{30} = 4.19\ \mathrm{mm}
  \end{align*}
  \hfill(2.0 points)

  Hence YES
  \hfill(1.0 points)
\end{description}

\solhead{7.15}{Flaring protoplanetary disk}
% source id: 2022_TH_5
\ioaafig{2022_TH_5-s1.pdf}
\begin{description}
\item[(a)] The angle between the light beam and the horizontal plane is
  \[
    \theta_l \approx \tan\theta_l = \frac{h(r)}{r}
  \]
  \hfill(1.0 points)

  Additionally, a tangent is inclined to the disc plane with an angle of
  \[
    \theta_T \approx \tan\theta_T = \frac{\Delta h(r)}{\Delta r}
  \]
  \hfill(2.0 points)

  An exterior angle of a triangle is equal to the sum of its two interior
  opposite angles. Thus,
  \[
    \beta = \theta_T - \theta_l = \frac{\Delta h(r)}{\Delta r} - \frac{h(r)}{r}
  \]
  \hfill(1.0 points)
\item[(b)] The flux from of stellar radiation at distance $r$ from the star is
  $E_s = \frac{L_s}{4\pi r^2}$, where $L_s$ is the luminosity of the star.
  However, the irradiation flux is the normal projection of this flux onto the
  infinitesimal small surface $A$ of the disk.
  \[
    Q_+ = E_s A\sin\beta \approx \frac{L_s A\beta}{4\pi r^2}
  \]
  \hfill(1.5 points)

  From the Stefan-Boltzmann law, the cooling rate is give by
  \[
    Q_- = A\sigma T_D^4.
  \]
  \hfill(0.5 points)

  In the thermal equilibrium
  \[
    Q_+ = Q_- \Longrightarrow
    T_D = \left(\frac{L_s\beta}{4\pi\sigma r^2}\right)^{\!\frac{1}{4}}
  \]
  \hfill(1.0 points)
\item[(c)] The condition for the isothermal layer can be solved by applying
  \begin{align*}
    \beta &= \frac{\Delta h(r)}{\Delta r} - \frac{h(r)}{r}
      = \frac{a(r+\Delta r)^b - ar^b}{\Delta r} - \frac{ar^b}{r} \\
    &= ar^b\left(\frac{(1+\frac{\Delta r}{r})^b - 1}{\Delta r} - \frac{1}{r}\right)
      = ar^b\left(\frac{b\Delta r}{r\Delta r} - \frac{1}{r}\right) \\
    &= a(b-1)r^{b-1} \propto r^2 \\
    &\Longrightarrow b = 3
  \end{align*}
  \hfill(2.0 points)

  Finally, one can determine the second constant by
  \begin{align*}
    \beta &= a(b-1)r^{b-1} = 2ar^2 \\
    \therefore\; T_{SL}^4 &= \left(\frac{aL_s}{2\pi\sigma}\right) \\
    a &= \left(\frac{2\pi\sigma T_{SL}^4}{L_s}\right)
  \end{align*}
  \hfill(1.0 points)
\end{description}

\solhead{7.16}{Neutrinos}
% source id: 2023_TH_Q8
\textbf{Neutrino emission of the supernova:}

From the table of constants, the core mass
$M = 1M_\odot = 1.988 \cdot 10^{30}$\,kg.
Mass of one atom of iron $^{56}_{26}\mathrm{Fe}$:
\[
  m_{\mathrm{Fe}} = 56\,\mathrm{Da}
  = 56 \times 1.661 \times 10^{-27}\,\mathrm{kg}
  = 9.3016 \times 10^{-26}\,\mathrm{kg}.
\]
(Note: atomic mass is expressed in \emph{atomic mass units}, symbol a.m.u.\ or
u, also called \emph{Daltons}, symbol Da.)

The number of $^{56}_{26}\mathrm{Fe}$ nuclei in the core is therefore:
\[
  n_{\mathrm{Fe}} = M/m_{\mathrm{Fe}} = 2.1373 \times 10^{55}
\]
\hfill(1 points)

and the initial number of nucleons is:
\[
  n_{\mathrm{nuc}} = 56\,n_{\mathrm{Fe}} = 1.1969 \times 10^{57}.
\]
\hfill(1 points)

Correspondingly, the initial number of protons is
$n_{\mathrm{p}} = 26\,n_{\mathrm{Fe}}$ and the initial number of neutrons is
$n_{\mathrm{n}} = 30\,n_{\mathrm{Fe}}$. \hfill(2 points)

We assume all the nucleons in the original $^{56}_{26}\mathrm{Fe}$ nuclei are
converted to individual nucleons in the given proportions
(1 proton : 8 neutrons). Therefore after the explosion,
\[
  n'_{\mathrm{n}} = \frac{8}{9} n_{\mathrm{nuc}} = 1.0634 \times 10^{57}
\]
neutrons and
\[
  n'_{\mathrm{p}} = \frac{1}{9} n_{\mathrm{nuc}} = 1.3299 \times 10^{56}
\]
protons remain, and thus
\[
  n_{\mathrm{p}} - n'_{\mathrm{p}} = 26\,n_{\mathrm{Fe}} - n'_{\mathrm{p}}
  = 4.2271 \times 10^{56}
\]
\hfill(4 points)

protons are changed. Since 1 neutrino is produced for every converted proton
($\mathrm{p} + \mathrm{e}^- \to \mathrm{n} + \nu_{\mathrm{e}}$), the neutrino
flux is:
\[
  I_\nu = \frac{4.2271 \times 10^{56}}{0.01\,\mathrm{s}}
  = 4.2271 \times 10^{58}\ \mathrm{neutrinos\ s^{-1}}.
\]
\hfill(2 points)

The flux density observed at Earth is given by:
\[
  F_{\mathrm{SN}} = \frac{I_\nu}{4\pi d_{\mathrm{Gal}}^2}.
\]
Substituting
$d_{\mathrm{Gal}} = 8\ \mathrm{kpc} = 8000 \times 3.086 \times 10^{16}\,\mathrm{m}
= 2.4688 \times 10^{20}\,\mathrm{m}$,
\[
  F_{\mathrm{SN}} \approx 5.5 \times 10^{16}\ \mathrm{neutrinos\ s^{-1}\,m^{-2}}.
\]
\hfill(2 points)

\textbf{Neutrino flux from the Sun:}

We can assume that the primary source of Solar luminosity is the p--p reaction,
\[
  4\mathrm{p}^+ + 2\mathrm{e}^- \to {}^4\mathrm{He}^{2+} + 2\nu_{\mathrm{e}}.
\]
Neglecting the electrons and neutrinos, the change in mass $\Delta m$ is:
\[
  \Delta m = 4 \cdot 1.673 \times 10^{-27}\,\mathrm{kg}
  - 6.645 \times 10^{-27}\,\mathrm{kg}
  = 4.7 \times 10^{-29}\,\mathrm{kg},
\]
and thus from $E = (\Delta m)c^2$ the amount of energy released is
${\approx}\,4.2244 \times 10^{-12}$\,J. Solar luminosity is
$L_\odot = 3.826 \times 10^{26}$\,W, so the number of reactions per second is
approximately:
\[
  \frac{3.826 \times 10^{26}}{4.2244 \times 10^{-12}} = 9.057 \times 10^{37}.
\]
With 2 neutrinos released per reaction, this gives a Solar neutrino flux of:
\[
  I_{\nu\odot} \approx 1.8 \times 10^{38}\ \mathrm{neutrinos\ s^{-1}}.
\]
(If the student remembers that the p--p reaction releases
$26.73\ \mathrm{MeV} = 4.2826 \times 10^{-12}$\,J, they will get
$8.9338 \times 10^{37}$ reactions per second and the same final answer.)
\hfill(getting to $I_{\nu\odot}$: 5 points)

At the distance of Earth, the observed flux density is given by:
\[
  F_\odot = \frac{I_{\nu\odot}}{4\pi d_\odot^2},
\]
where $d_\odot = 1\ \mathrm{au} = 1.496 \times 10^{11}$\,m, and so:
\[
  F_\odot \approx 6.4 \times 10^{14}\ \mathrm{neutrinos\ s^{-1}\,m^{-2}}.
\]
\hfill(1 point)

The final ratio is therefore:
\[
  F_{SN}/F_\odot \approx \frac{5.5 \times 10^{16}}{6.4 \times 10^{14}}
  \approx 86 \approx 100,
\]
or two orders of magnitude. \hfill(2 points)

\solhead{7.17}{Stars through graphs}
% source id: 2025_TH_T11
\begin{description}
\item[(T11.1a)]\mbox{}
  \begin{itemize}
  \item $X$ and $l$ increase with $r$, while $T$, $\epsilon_{\mathrm{nuc}}$,
    and $Y$ decrease with $r$. So A and B must be $X$ and $l$.
  \item For a $1\,\mathrm{M_\odot}$ star of age 4\,Gyr, $X \neq 0$ at centre.
    So B cannot be $X$.
  \item Also, $l = 0$ at centre. Matches with B.
  \item Both $X$ and $l$ must saturate to their maximum values just outside
    the H-burning core, and this happens at $r/R = 0.3$ here.
  \item Each of $T$, $\epsilon_{\mathrm{nuc}}$, and $Y$ have their maximum
    value at the centre and decrease outwards. So C is one of these.
  \item But $Y$ cannot be (close to) zero at the surface ($r/R = 1$) compared
    to its central value.
  \item $\epsilon_{\mathrm{nuc}}$ becomes extremely small just outside the
    core ($r/R \lesssim 0.3$), and not at the surface.
  \item Therefore, C must be $T(r)$ (the surface value is
    ${\sim}\,10^3$\,K, compared to the maximum value at the centre
    ${\sim}\,10^6$\,K).
  \end{itemize}
  Therefore,
  \begin{align*}
    \mathrm{A} &\longrightarrow X(r) \tag*{(2.0)}\\
    \mathrm{B} &\longrightarrow l(r) \tag*{(2.0)}\\
    \mathrm{C} &\longrightarrow T(r) \tag*{(2.0)}
  \end{align*}

\item[(T11.1b)] Maximum value of $X$ occurs outside the core, and is given by
  \[
    X_{\mathrm{s}} = 1 - Y_{\mathrm{s}} - Z_{\mathrm{s}}
    = 1 - 0.28 - 0.02 = 0.70
  \]
  \hfill(1.0)

  Note $Z$ remains constant throughout the star in the absence of
  gravitational settling and diffusion (given).

  From the given graph, the central hydrogen mass fraction, $X_{\mathrm{c}}$,
  is approximately
  \[
    X_{\mathrm{c}} \approx 0.55 X_{\mathrm{s}} = 0.55 \times 0.70 = 0.39
  \]
  \hfill(1.5)

  Then, the helium mass fraction at the centre is
  \[
    Y_{\mathrm{c}} = 1 - X_{\mathrm{c}} - Z_{\mathrm{c}}
    = 1 - 0.39 - 0.02 = 0.59
  \]
  \hfill(0.5)

\item[(T11.1c)]\mbox{}
  \begin{itemize}
  \item The curve for $\epsilon_{\mathrm{nuc}}$ must mirror that of $l$: any
    reasonable smooth curve between $(0,1)$ and $(0.3,0)$. \hfill(2.5)
  \item The curve for $Y$ must mirror that of $X$, and flatten outside the
    core to the normalized value corresponding to $Y_{\mathrm{s}} = 0.28$.
    Previous result of $Y_{\mathrm{c}} = 0.59$ gives this normalized value:
    $0.28/0.59 \approx 0.47$. \hfill(2.5)
  \end{itemize}
  \ioaafig{2025_TH_T11-s1.pdf}

\item[(T11.2a)] During the main sequence phase, the mass fraction of helium at
  the centre increases steadily, up to a maximum value of
  $Y_{\mathrm{c,\,max}} = 1 - Z_{\mathrm{c}}$. Here that maximum value is
  reached at an age of around 8.5\,Gyr, from the plot of $Y_{\mathrm{c}}$ vs
  age.
  \[
    t_{\mathrm{MS}} = 8.5 \times 10^{9}\,\mathrm{yr}
  \]
  \hfill(1.0)

\item[(T11.2b)] During the helium burning phase, the mass fraction of helium
  at the centre decreases steadily, down to 0. The inset in the plot of
  $Y_{\mathrm{c}}$ vs age shows that this corresponds to ages between
  $11.896 \times 10^{9}$\,yr and $11.960 \times 10^{9}$\,yr, corresponding to
  a duration of 64\,Myr.
  \[
    \Delta t_{\mathrm{He}} = 64 \times 10^{6}\,\mathrm{yr}
  \]
  \hfill(1.0)

\item[(T11.2c)] Read-off from the plot of $\log L/L_\odot$ vs age, we get age
  ${\approx}\,4 \times 10^{9}$\,yr. \hfill(1.0)

  At this age, read-off from plot of $Y_{\mathrm{c}}$ vs age, we get
  \[
    Y_{\mathrm{c}} \approx 0.60 \implies
    X_{\mathrm{c}} \approx 1 - 0.60 - 0.02 = 0.38
  \]
  \hfill(1.0)

  Therefore the fraction of original H burnt is
  \[
    f_{\mathrm{H}} = 1 - \frac{X_{\mathrm{c}}}{X_0}
    = 1 - \frac{X_{\mathrm{c}}}{1 - Y_0 - Z_0} = 0.46
  \]
  \hfill(1.0)

\item[(T11.2d)]
  \[
    X_0 = 1 - Y_0 - Z_0 = 1 - 0.28 - 0.02 = 0.70
  \]
  At this stage 60\% of H has been burnt at the centre. So
  $X_{\mathrm{c}} = 0.40 X_0 = 0.28$. Therefore,
  \[
    Y_{\mathrm{c}} = 1 - X_{\mathrm{c}} - Z_{\mathrm{c}}
    = 1 - 0.28 - 0.02 = 0.70
  \]
  \hfill(2.0)

  From the graph of $Y_{\mathrm{c}}$ vs age, this corresponds to an age of
  approximately $5 \times 10^{9}$\,yr. At this age, the value of radius is
  approximately $1\,\mathrm{R_\odot}$ ($\log R/R_\odot \approx 0$), from the
  graph of $\log R/R_\odot$ vs age.
  \[
    R_1 = 1\,\mathrm{R_\odot}
  \]
  \hfill(1.0)

\item[(T11.2e)] For P:
  $\log T_{\mathrm{eff}} = 3.74 \implies T_{\mathrm{eff}} \approx 5500$\,K
  (the $L$, being flat, does not help here). \hfill(1.0)

  Read-off from plot of $T_{\mathrm{eff}}$ vs age, we get age
  ${\approx}\,11 \times 10^{9}$\,yr. At this age, read-off from plot of
  $\log R/R_\odot$ vs age, we get $\log R/R_\odot \approx 0.25 \implies
  R_{\mathrm{P}} = 1.78\,\mathrm{R_\odot}$. \hfill(1.0)

  For Q: $\log L/L_\odot = 2.5$. (One can also use
  $\log T_{\mathrm{eff}} = 3.60 \implies T_{\mathrm{eff}} \approx 3981$\,K
  which would require rounding off to 3950\,K at the next step.) \hfill(1.0)

  Read-off from (the inset) plot of $\log L/L_\odot$ vs age, we get age
  ${\approx}\,11.884 \times 10^{9}$\,yr. Notice that this value of
  $\log L/L_\odot$ occurs twice, and one has to take the first one, i.e., the
  ascending RGB branch.

  At this age, read-off from (inset) plot of $\log R/R_\odot$ vs age, we get
  $\log R/R_\odot \approx 1.6 \implies R_{\mathrm{Q}} = 39.81\,\mathrm{R_\odot}$.
  \hfill(1.0)

\item[(T11.3a)]
  \begin{align*}
    \frac{dq}{dx} &= \frac{dq}{dm}\frac{dm}{dr}\frac{dr}{dx}
      \tag*{(0.5)}\\
      &= \frac{1}{M}\left[4\pi (Rx)^2 \rho\right] R\\
      &= 3\left(\frac{4\pi R^3}{3M}\right) x^2 \rho
       = 3 \frac{1}{\bar{\rho}} x^2 \rho
      \tag*{(1.0)}\\
    \therefore\ \frac{dq}{dx} &= 3 x^2 \sigma
      \tag*{(0.5)}
  \end{align*}

\item[(T11.3b)] Let us denote the line as
  \[
    \log\sigma = \alpha \log x + \beta
  \]
  From the graph, we use the two end points $(-0.5, 1.0)$ and $(-0.1, -1.0)$.
  Therefore,
  \begin{align*}
    1.0 &= -0.5\alpha + \beta\\
    -1.0 &= -0.1\alpha + \beta
  \end{align*}
  Solving, we get $\alpha = -5$ and $\beta = -1.5$. Therefore,
  \begin{align*}
    \log\sigma &= -5\log x - 1.5
      &&\text{for } -0.5 < \log x < -0.1
      \tag*{(2.0)}\\
    \therefore\ \sigma(x) &= 10^{-1.5} x^{-5}
      &&\text{for } 0.32 < x < 0.80\\
    \implies \sigma(x) &= 0.030\,x^{-5}
      &&\text{for } 0.32 < x < 0.80
      \tag*{(2.0)}
  \end{align*}

\item[(T11.3c)]
  \[
    \frac{dq}{dx} = 3x^2\sigma = 3x^2(0.030\,x^{-5}) = 0.090\,x^{-3}
  \]
  \hfill(1.0)

  Integrating, we get
  \begin{align*}
    q(x) &= \int 0.090\,x^{-3}\,dx
      \tag*{(1.0)}\\
      &= -\frac{0.045}{x^2} + C
      \qquad (C \text{ is a constant of integration})
      \tag*{(1.0)}
  \end{align*}
  Now, given that $q(x = 0.35) = 0.7$. Using this, we get $C = 1.067$.
  \hfill(2.0)

  Therefore (rounding off to within 0.005),
  \[
    q(x) = 1.065 - \frac{0.045}{x^2}
    \qquad \text{for } 0.32 < x < 0.80
  \]
  \hfill(1.0)

\item[(T11.3d)]
  \[
    \frac{dq}{dx} = 3x^2\sigma = 3x^2(Ax + B) = 3Ax^3 + 3Bx^2
  \]
  \hfill(1.0)

  Integrating, we get
  \begin{align*}
    q(x) &= \int (3Ax^3 + 3Bx^2)\,dx
      \tag*{(1.0)}\\
      &= \frac{3A}{4}x^4 + Bx^3 + D
      \qquad (D \text{ is a constant of integration})
      \tag*{(1.0)}
  \end{align*}
  Three boundary conditions applied sequentially to find the constants:
  \begin{itemize}
  \item Mass at the centre is zero, i.e., $q(0) = 0 \implies D = 0$
    \hfill(0.5)
  \item $\rho(0) = 80\bar{\rho}$ (given), i.e.,
    $\sigma(0) = 80 \implies B = 80$ \hfill(1.5)
  \item Given $q(0.25) = 0.50$. Therefore,
    \begin{align*}
      0.5 &= \frac{3A}{4}(0.25)^4 + 80\,(0.25)^3\\
      \implies A &= \frac{4}{3}\,\frac{0.50 - 80(0.25)^3}{(0.25)^4}\\
      \implies A &= -256
      \tag*{(2.0)}
    \end{align*}
  \end{itemize}
  Therefore, putting the values of $A$ and $B$,
  \[
    q(x) = -192 x^4 + 80 x^3 \qquad \text{for } 0 \le x \le 0.32
  \]
  \hfill(1.0)

\item[(T11.3e)] The sketch should be a smooth curve for $q(x)$ from $(0,0)$
  to $(1,1)$, with zero slope at both ends, passing through the two given
  points $(0.25, 0.5)$ and $(0.35, 0.7)$; the initial part looks like a
  cubic--quartic, and the latter part looks like a negative inverse-square
  with a dc shift.
  \ioaafig{2025_TH_T11-s2.pdf}
\end{description}

